\title{LongRoPE: Extending LLM Context Window Beyond 2 Million Tokens}

\begin{abstract}
Expanding the context window in large language models (LLMs) is highly sought after, yet it remains restricted—typically up to 128k tokens—because fine-tuning such models is expensive, lengthy data is in short supply, and additional token positions often result in problematic values.

This study presents LongRoPE, which, for the first time, scales the context window of pre-trained LLMs up to an impressive \textbf{2048k} tokens, requiring as few as 1k fine-tuning steps at training lengths of 256k, all while preserving the model’s original short-context window capabilities. This advancement is grounded in three central techniques: (i) through efficient searching, we detect and utilize two types of non-uniformities present in positional interpolation, enhancing fine-tuning initialization and enabling an eightfold context expansion even without fine-tuning; (ii) we propose a staged extension procedure, initially fine-tuning the model on 256k token contexts before applying a subsequent round of positional interpolation to the already fine-tuned model, resulting in a 2048k context window; and (iii) we recalibrate LongRoPE on 8k token contexts to restore short-window performance. Comprehensive empirical evaluation on LLaMA2 and Mistral across a spectrum of tasks confirms the method's efficacy. Importantly, models enhanced through LongRoPE keep their original architecture—making only slight adjustments to the positional embedding—and are compatible with most established optimizations. Code will be released at \texttt{https://github.com/microsoft/LongRoPE}

\section{Introduction}

Although Large Language Models (LLMs) have demonstrated outstanding capabilities across a range of applications~\cite{gpt4,llama2}, they are constrained by relatively short context window lengths; for instance, LLaMA2 imposes a restriction of 4096 tokens~\cite{llama2}. When input sequences exceed this threshold, model performance degrades, as the model has not been exposed to such extended positional contexts during training. This limitation creates significant difficulties in key areas, such as in-context learning that requires processing many examples~\cite{Huang2023FewerIM}, as well as for LLM-driven agents~\cite{park2023generative,madaan2023selfrefine}.

Recent studies demonstrate that by fine-tuning on longer sequences, a pre-trained LLM's context window can reach up to approximately 128k tokens~\cite{longlora,pi,yarn,zhang2024soaring,liu2023scaling}. However, advancing this limit presents three main challenges. First, adding untrained position indices brings in numerous unstable values, resulting in out-of-distribution complications that hinder the convergence of fine-tuning~\cite{pi}. This problem intensifies when scaling from 4k to over 1000k, since the majority ($>$90\%) of positions are novel. Second, fine-tuning generally depends on access to texts of matching lengths, yet datasets featuring such lengthy samples—especially beyond 1000k tokens—are scarce. Furthermore, processing extremely long sequences during training demands substantial computational resources, including excessive time and GPU usage. Third, expanding to ultra-large context windows leads to diluted attention, as it is distributed across a vast array of token positions, which undermines performance on tasks involving shorter contexts~\cite{pi}.

To address the initial issue, a commonly employed strategy is to interpolate RoPE positional embeddings~\cite{rope,pi}, which involves mapping new positional indices to fit within the scope utilized during pretraining, as depicted in Fig.\ref{fig:overview}. Position Interpolation (PI)~\cite{pi} operates by applying a linear interpolation to the rotary angles in RoPE according to the expansion ratio. NTK~\cite{ntk,dynamicntk}, in contrast, proposes the use of non-uniform interpolation and extrapolation across different dimensions of RoPE. YaRN~\cite{yarn} takes a frequency-based approach, segmenting RoPE dimensions into three groups and performing extrapolation, NTK-based interpolation, and linear interpolation for each, respectively. Despite these methods, positional embeddings within the Transformer architecture inherently possess a \emph{complex, non-uniform distribution of information entropy}. Current techniques fail to adequately exploit this subtle heterogeneity, resulting in the loss of valuable information and thereby restricting the attainable context window size.

Section~\ref{sec:analysis} presents two principal empirical insights: \textbf{(1)} Effective positional interpolation must account for two distinct types of non-uniformity—those arising from the distribution across RoPE dimensions and those due to variations in token positions. In particular, both the lower dimensions of RoPE and tokens at the sequence's beginning require less interpolation, although the precise optimal approach is contingent upon the length to which the model is extended. 
\textbf{(2)} Integrating an awareness of these non-uniformities into the positional interpolation process significantly preserves the original RoPE information, especially in crucial dimensions and at key token positions. This approach reduces interpolation-induced information loss, thereby improving initial conditions for subsequent fine-tuning. Furthermore, it makes it possible to achieve up to an 8$\times$ extension in contexts that do not involve fine-tuning.

Driven by these observations, we propose \textbf{LongRoPE}, a robust approach for increasing the LLM context window to more than 2 \emph{million} tokens. LongRoPE leverages three principal advances. Firstly, LongRoPE maximally utilizes multidimensional variations in positional interpolation, pinpointing optimal rescale factors for RoPE rotation angles on a per-dimension basis according to token position. Since the search space for discovering appropriate rescale factors grows exponentially with the desired extension ratio, LongRoPE implements an evolutionary search procedure enhanced by two optimization strategies to improve the efficiency of this process. Fig.~\ref{fig:overview} illustrates a representative example of the resulting rescaled RoPE.

Subsequently, {LongRoPE} employs a resourceful, stepwise extension method, enabling a 2048k context window without requiring direct fine-tuning on extremely long textual inputs, which are both rare and difficult to obtain. The process initiates by identifying a 256k sequence length suitable for the pre-trained LLM and then fine-tuning the model to operate at this length. Since our non-uniform positional interpolation facilitates an 8$\times$ expansion without additional fine-tuning, a subsequent search is performed to determine appropriate RoPE rescale factors for the newly fine-tuned and length-extended LLM. This approach results in the successful attainment of a 2048k context window for models like LLaMA2 and Mistral~\cite{mistral}.

Lastly, to prevent reduced effectiveness on the original, smaller context window, {LongRoPE} maintains adaptation of the RoPE rescale factors even after extending the LLM. Analogous to the process of scaling up from 256k to 2048k, our method also applies the search algorithm to downscale to 4k and 8k context windows on the LLM fine-tuned at 256k, thereby minimizing the reliance on positional interpolation. At inference time, if the input sequence is under 8k tokens in length, RoPE is recalibrated with the optimized rescale factors identified through this search.

A wide range of experiments on multiple LLMs and diverse long-context evaluation tasks confirm the efficacy of our proposed approach. Our findings indicate that LongRoPE excels at preserving low perplexity across evaluation lengths spanning from 4k up to 2048k, surpasses 90\% accuracy in passkey retrieval, and matches standard benchmark performance constructed for the 4096-token context window. The LongRoPE technique is universally compatible with any LLM utilizing RoPE embeddings. Both our code and the LongRoPE-2048k model checkpoints will be made publicly available.

\section{Non-uniformity in Positional Interpolation}
\label{sec:analysis}

\subsection{Preliminary}

Transformer models require explicit positional information, often in the form of position embedding, to represent the order of input tokens. Our work focuses on the RoPE~\cite{rope} position embedding, which is widely used in recent  LLMs. For a token at position index $n$, its corresponding RoPE encoding can be simplified as follows:

\begin{equation}
		\fontsize{7.8}{7.8} \selectfont
	\label{eq:rope}
	\left[ cos(n\theta_0), sin(n\theta_0), cos(n\theta_1), \cdot\cdot\cdot, cos(n\theta_{d/2-1}), sin(n\theta_{d/2-1}) \right]   
\end{equation}

Here, $d$ denotes the dimension of the embedding, $n\theta_i$ refers to the rotary phase assigned to a token located at position $n$, and $\theta_i=\theta^{-2i/d}$ characterizes the frequency of rotation. The standard value for $\theta$ in RoPE is set at 10000.

\textbf{\emph{Context window extension ratio $s$} and positional interpolation}. The extension ratio $s$ is defined as the proportion between the new context length $L'$ and the original context length $L$, i.e., $s=\frac{L'}{L}$.

To extend the context window from $L$ to $L'$, current positional interpolation methods suggest downscaling rotation frequencies $\theta_i$ based on extension ratio $s$. Let $\beta=\theta^{2/d}$, and $\lambda$ denote the actual rescale factor related to $s$, we   unify these positional interpolation methods as follows:

\begin{equation}	
	\fontsize{7.5}{7.5} \selectfont
	\label{eq:unified_rope}
	\begin{aligned}
		\left[cos\left(\frac{n}{\lambda(\beta)^0}\right), sin \left(\frac{n}{\lambda(\beta)^0}\right), cos\left(\frac{n}{\lambda(\beta)^1}\right), \cdot\cdot\cdot,  sin \left(\frac{n}{\lambda(\beta)^{d/2-1}}\right) \right]   
	\end{aligned}
\end{equation}

\textbf{Linear positional interpolation (PI)}. PI~\cite{pi} proposes the application of linear interpolation to position indices that fall within the original pre-trained sequence length. When extending the context by a factor of $s$, all positional rotation angles are uniformly scaled down by $\lambda = s$ in every RoPE dimension. This approach, however, results in a compression of positional representations, making position encodings less distinguishable for tokens situated near each other. Consequently, PI's performance typically degrades as the extension ratio increases.

\textbf{NTK-based interpolation and extrapolation}. ~\cite{ntk,dynamicntk} approach RoPE from the standpoint of information encoding, employing Neural Tangent Kernel (NTK) theory~\cite{jacot2018neural,tancik2020fourier} in their analysis. In order to address the issue of congestion at specific positions in PI, these works propose spreading the interpolation effect more evenly across the various dimensions of RoPE. Specifically, the scaling applied is lesser for the lower (high frequency) dimensions and greater for the higher (low frequency) ones, which facilitates both interpolation and extrapolation of positional embeddings, with $\lambda=s^{i}$. The enhanced dynamic NTK~\cite{dynamicntk} method further refines this approach by allowing the extension ratio to be position-dependent and responsive to the length of the input sequence. While PI requires further fine-tuning when extending context lengths, NTK-aware approaches can enlarge context windows even in scenarios where fine-tuning is not feasible, though in practice the window extension is typically limited to a 4$\times$ ratio.

\textbf{YaRN}~\cite{yarn} presents a substantial advancement in the efficacy of positional interpolation. The method organizes RoPE dimensions into three distinct groups based on their frequency characteristics, applying a unique interpolation technique to each group. Extrapolation with $\lambda=1$ is used for the highest frequency dimensions, linear interpolation (PI) is reserved for those with the lowest frequencies, and the NTK approach is applied to intermediate frequencies. Central to YaRN's methodology is this frequency-based partitioning of RoPE dimensions, which is presently determined through manual empirical processes. As a result, this approach may not yield optimal outcomes when applied to novel LLMs.

\subsection{Study on Non-uniform Positional Interpolation}

Drawing motivation from NTK and YaRN, we observe that their improvements stem from introducing non-linearity—particularly by treating distinct frequencies across RoPE dimensions uniquely, enabling targeted interpolation and extrapolation. Nonetheless, prevailing non-linear approaches predominantly depend on handcrafted human heuristics. This observation leads to two primary inquiries: (1) Does the existing positional interpolation approach represent the optimum? (2) Might there exist alternative, yet uninvestigated, forms of non-linearity?

In pursuit of these answers, we employ evolution search (refer to Sec.~\ref{sec:method}) to identify improved non-uniform positional interpolations for LLaMA2-7B. This search process is steered by perplexity measurements, utilizing five randomly chosen examples from the PG19~\cite{pg19} validation dataset. Our experimental evaluation uncovers several notable conclusions.

\textbf{Finding 1}: \textit{RoPE dimensions exhibit substantial non-uniformities, which are not effectively handled by current positional interpolation methods.}

We determine the optimal value of $\lambda$ for each RoPE dimension as defined in Eq.~\ref{eq:unified_rope}. Table~\ref{tbl:exp1} reports the perplexity scores of LLaMA2-7B using different approaches on the PG19 and Proof-pile~\cite{proof-pile} test sets, evaluated without additional fine-tuning. The solution derived from our search yields marked perplexity reductions, indicating that traditional linear (PI) as well as previously proposed non-uniform interpolation schemes (Dynamic-NTK and YaRN) fail to realize optimal performance. Interestingly, YaRN is less effective than PI and NTK on the PG19 benchmark because, for language models that have not been fine-tuned, it fails to attain the target length of the context window. Specifically, the perplexity associated with YaRN increases sharply beyond 7,000 tokens in an 8,000-token context.

Our exploration reveals that the scaling factors $\lambda$ in Eq.~\ref{eq:unified_rope} are heterogeneous rather than constant, setting them apart from the uniform scaling $s$ applied in PI, NTK’s direct formula, or the group-level adjustments in YaRN. Employing these non-uniform scaling factors yields a pronounced enhancement in LLaMA2’s language modeling outcomes, particularly regarding perplexity at context lengths of 8k and 16k, and crucially, this is achieved without additional fine-tuning. The efficacy of this approach stems from the revised positional embeddings' capacity to closely approximate the original RoPE—particularly preserving essential dimensions—which in turn alleviates the model's challenge in distinguishing between tokens that are nearby in sequence.

\textbf{Finding 2}: \textit{RoPE for the initial tokens in the input sequence should be extrapolated with less interpolation.} 

We propose that for the first $\hat{n}$ tokens in each input sequence, RoPE should apply minimal or no interpolation. This rationale stems from their elevated attention scores, as documented in Streaming LLM~\cite{streamingllm} and LM-Infinite~\cite{han2023lminfinite}, indicating their essential role in attention mechanisms. To test this premise, we increase the context window size to 8k and 16k with PI and NTK schemes, but maintain the interpolation-free processing for the initial $\hat{n}$ tokens ($0, 2, ..., 256$). Setting $\hat{n}=0$ recovers the conventional PI and NTK approach. Table~\ref{tbl:tokenposition} presents two main findings: \textbf{(1)} preserving the position representation of early tokens without interpolation produces performance gains, and \textbf{(2)} the optimal value of $\hat n$ varies based on the intended context window expansion.

\textbf{Finding 3}: \textit{Non-uniform positional interpolation effectively extends LLM context window in both fine-tuning and non-fine-tuning settings.} 

Although our experiments have demonstrated that the non-uniform position interpolation we identified yields substantial gains in extension capability at 8k and 16k contexts without any fine-tuning, the expansion to even longer contexts necessitates further fine-tuning. Consequently, we conduct fine-tuning of LLaMA2-7B using our optimized RoPE approach for a 64k context window (refer to Appendix for experimental configurations). As indicated in Table~\ref{tbl:finetune}, our approach consistently surpasses both PI and YaRN, outperforming them on LLaMA2-7B both prior to and following the fine-tuning process. This improvement stems from the strategic application of non-uniform positional interpolation, which reduces information degradation and enables superior initialization during fine-tuning.

\textbf{Summary}. The research identifies two sources of non-uniformity: disparities among RoPE dimensions and variations in token positions. Harnessing these non-uniform patterns in the context of positional interpolation delivers significant enhancements in the ability of LLMs to operate effectively with extended context windows.

\section{LongRoPE}

\label{sec:method}
Inspired by these observations, we propose LongRoPE. This method initially employs an optimized search algorithm designed to harness both forms of non-uniformity, and subsequently applies this approach to push the LLM context window past 2 million tokens.

\subsection{Problem Formulation}

These two forms of non-uniformity significantly expand the solution space and complicate the optimization process. To manage these challenges, we recast the optimization task of multidimensional non-uniform position interpolation as a search problem.

For a LLM targeting a  context window size of $L'$ and lengthy input documents $\mathbf{X}$, where each $\mathbf{x}\in\mathbf{X}$ surpasses $L'$ in token length, we denote the original rotary angle of the $i^{th}$ dimension in RoPE embedding at token position $n$ as  $\frac{n}{\beta^i}$. The optimization problem is then formulated as follows:

\begin{equation}
\begin{gathered}
		\label{eq:problem}

\underset{\mathbf{x}\in\mathbf{X};\, |\mathbf{x}|\ge L'}{\operatorname{arg\,min}}\, \mathcal{L}\left(\text{LLM}(\text{RoPE}, \mathbf{X})\right),\, \text{with}
\\
\scalebox{0.93}{
$\underset {\begin{subarray}{l} i=0,\dotsc,\frac{d}{2}-1;\\ n\in[0,|\mathbf{x}|);\end{subarray}}{\text{RoPE}_{(n)}} = {\left[\dotsc, \cos\left( \mathbb{I}(\hat{\lambda}_i, \hat{n}) \frac{n}{\beta^i} \right),\, \sin\left( \mathbb{I}(\hat{\lambda}_i, \hat{n}) \frac{n}{\beta^i} \right),\, \dotsc \right]}$}
\\
\text{where}\; \mathbb{I}(\hat{\lambda}_i, \hat{n}) = \begin{cases}
1 & \text{if}\; n < \hat{n} \\
{\frac{1}{\lambda}_i} & \text{if}\; n \geq \hat{n}
\end{cases}
\end{gathered}
\end{equation}

We introduce rescale factors, $\mathbb{I}(\hat{\lambda}_{i},\hat{n})$, to address two types of non-uniformity: $\hat{\lambda_i}$ characterizes non-uniformity across RoPE dimensions, while $\hat{n}$ reflects positional irregularities among tokens. Specifically, $\mathbb{I}(\hat{\lambda}_{i},\hat{n})$ is employed to adjust the rotation angle for the $i^{th}$ dimension of RoPE, with $\hat\lambda_{i}$ as its corresponding rescale value and $\hat{n}$ serving as the threshold for token positions. For tokens at positions less than $\hat{n}$, i.e., the first $\hat{n}-1$ tokens, the rescale factor $\hat\lambda_{i}$ remains inactive and the standard RoPE rotational angle $\frac{n}{\beta^i}$ is utilized. When token positions reach or exceed $n\ge\hat{n}$, the rescale factor is activated and incorporated.

For a desired context window length of $L'$, the task is to determine the set of optimal rescaling parameters ($\mathbb{I}(\hat{\lambda}_{0},\hat{n})$, $\mathbb{I}(\hat{\lambda}_{1},\hat{n})$, ... $\mathbb{I}(\hat{\lambda}_{i},\hat{n})$, ...) for each RoPE dimension from the $1^{st}$ up to the $d^{th}$. Implementing these rescale factors within the $\text{RoPE}$ module of the target $\text{LLM}$ enables the model to minimize the next-token prediction loss, denoted by $\mathcal{L}$ (corresponding to perplexity), when processing input sequences $\mathbf{X}$ whose lengths are $L'$.

\subsection{Searching the Non-uniform Position Interpolation}

In order to address the challenge formulated in Eq.~\ref{eq:problem}, we propose a straightforward but powerful approach that identifies the optimal RoPE rescaling coefficients. This strategy is specifically designed to take advantage of the diverse, multidimensional variances present in position embeddings.

\textbf{Search space}. We establish an extensive search space to accommodate both types of non-uniformity. Table~\ref{tbl:searchspace} provides an overview of the search space configuration. In detail, our approach enables exploration of tailored rescale factors assigned individually to each dimension within the RoPE embedding. To streamline the construction of the search space, we optimize over $\lambda_i$ and $\hat{n}$, rather than directly tuning $\mathbb{I}(\hat{\lambda}_{i},\hat{n})$; in this context, $\hat{\lambda}_{i}$ is defined as $1/\lambda_i$. As indicated in Table~\ref{tbl:searchspace}, the allowable range for $\lambda_i$ spans from a minimum of 1.0, corresponding to simple extrapolation, up to a maximum of $s\times1.25$, which allows for greater interpolation compared to PI. The increment for $\lambda_i$ during the search is set to 0.01, with $s$ denoting the intended extension ratio for the context window.

The parameter $\hat{n}$ determines how many of the earliest token positions maintain the original RoPE embedding rather than undergoing position interpolation. In our experiments, we vary $\hat{n}$ over the set \{0, 1, 2, 4, 8, 12, 16, 20, 24, 28, 32, 64, 128, 256\} to find optimal values. Setting $\hat{n}=0$ implies that every token position employs the optimized rescale factors throughout.

\textbf{Evolution-based search}. The search space detailed in Table~\ref{tbl:searchspace} encompasses a wide array of positional interpolation options, making exhaustive exploration highly non-trivial. For instance, scaling to a $s=4\times$ extension yields a total of $400^{128/2}\times14 = 4\times10^{167}$ possibilities. As the extension ratio increases, the size of the search space grows at an exponential rate. To mitigate this complexity, we employ evolution search~\cite{guo2020single}, and incorporate two optimization strategies that significantly enhance the efficiency of the search. Algorithm~\ref{alg:search} presents an overview of the search workflow.

\textit{Enhanced initial population creation}. Rather than generating $P$ rescale factors at random for the initial population, we incorporate the three RoPE rescale factors associated with PI, NTK, and YaRN as distinct members of this population. For the other $P$-3 individuals, we apply random mutations to the three rescale factors, each with a mutation likelihood of $p$.

\textit{Monotonically non-decreasing constraint}. Once the initial population is constructed, we evaluate LLM perplexity for every member. To do this, the associated RoPE rescale factors are implemented in the target LLM, followed by perplexity calculation on the input $\mathbf{X}$. The top-k candidates, as determined by their performance, are selected as parents for the evolutionary process. Due to the expansive nature of the search space, simple mutation and crossover may frequently yield suboptimal candidates, thus necessitating redundant perplexity evaluations. This inefficiency is magnified when $L'$ is substantial, since each perplexity assessment requires considerable computational effort.

To mitigate this issue, we enforce a monotonicity constraint such that the rescaled RoPE factors are non-decreasing: $\lambda_i\le \lambda_{i+1}$. Only those RoPE settings that conform to this constraint are utilized in the LLM perplexity assessment, which markedly cuts down the computational expense involved in searching. Concretely, we stipulate that $\lambda_i$ should follow an increasing trend across RoPE dimensions (namely, for $i$=0,...,63). This constraint reflects insights from NTK theory~\cite{jacot2018neural,tancik2020fourier,ntk}, which posits that dimensions corresponding to higher frequency (lower indices) require less interpolation, implying a smaller value for $\lambda_i$, while those associated with lower frequency (higher indices) benefit from greater interpolation, permitting a larger $\lambda_i$.

\begin{algorithm}[t]
	\small
	\caption{The search algorithm}
	\textbf{Input:} target LLM, input samples $\mathbf{X}$, population size $P$, mutation size $N_1$, crossover size $N_2$, max iterations $\mathcal{T}$, mutation probability $p$\\

	\begin{algorithmic}[1]
		\label{alg:search}
		\STATE Top-k $\leftarrow$ $\phi$;	
		\STATE $\text{P}_0$ $\leftarrow$ \textit{Initialize\_population\_with\_optimization} ($P$, $\mathbf{X}$, $p$); 
		\FOR{$i=1$ to $\mathcal{T}$}
		\STATE \textit{Compute\_perplexity} (LLM, $\text{P}_{i-1}$, $\mathbf{X}$);
		\STATE Top-k $\leftarrow$ \textit{Update\_Topk} (Top-k, $\text{P}_{i-1}$);
		\STATE $\text{P}_{mutation}$ $\leftarrow$ \textit{Mutation\_with\_mono\_constraint} (Top-k, $N_1$, $p$);
		\STATE $\text{P}_{crossover}$ $\leftarrow$ \textit{Crossover\_with\_mono\_constraint} (Top-k, $N_2$);
		\STATE $\text{P}_i$ = $\text{P}_{mutation}$ $\cup$ $\text{P}_{crossover}$ $\cup$ Top-k;
		\ENDFOR
		\STATE Output the individual in Top-k exhibiting the lowest perplexity;
	\end{algorithmic}
\end{algorithm}

\textbf{8$\times$ extension without fine-tuning}. Through evolutionary search, our approach efficiently uncovers non-uniform RoPE rescale factors that selectively retain crucial dimensions and positions, thereby reducing the loss of information caused by interpolation. As illustrated in Fig.\ref{fig:non-ft}, this enables us to expand LLaMA2’s context window from 4k up to 32k tokens without the need for any fine-tuning. By comparison, other techniques like PI as well as non-uniform NTK and YaRN experience significant increases in perplexity following a 2$\times$ context window extension.

\subsection{Extending LLM Context Window to 2048K}
\label{sec:extendto2048k}

\textbf{Progressive extension to 2048k}. In this section, we present an approach for expanding the context window in pre-trained LLMs from the conventional 4k up to more than 2048k. Our non-uniform positional interpolation strategy permits up to an 8$\times$ increase without requiring any fine-tuning. When much greater extensions, such as 512$\times$, are necessary, fine-tuning becomes indispensable. A typical strategy involves identifying suitable RoPE rescaling factors corresponding to the desired 2048k context length followed by fine-tuning. Nonetheless, this technique is confronted by significant obstacles, namely the immense computational cost involved. Furthermore, according to our practical experience, achieving effective fine-tuning at such high extension ratios is notably difficult (refer to Appendix).

Fortunately, LongRoPE demonstrates efficacy in both pre-trained and further fine-tuned extended LLMs. As a result, we present a streamlined, incremental approach that reaches the desired 2048k context window using only 1k fine-tuning iterations, all within a training sequence length capped at 256k.

$\Diamond$ \textit{Scaling pre-trained LLM context windows to 256k with LongRoPE search}. Using LLaMA2 for illustration, we perform a search targeting context window lengths of 128k and 256k tokens. The resulting context extension factors achieved in this step are 32$\times$ and 64$\times$, respectively.

$\Diamond$ \textit{Fine-tuning to 256k}. Subsequently, we adapt the pre-trained LLM to accommodate a 256k context window. More concretely, LLaMA2 is initially fine-tuned for 400 steps utilizing RoPE rescaling tailored to 128k. After this, we substitute the RoPE scaling factors with those corresponding to 256k on the already fine-tuned checkpoint, and proceed with a further 600 steps of fine-tuning. This two-stage process demonstrates greater efficiency compared to immediate fine-tuning at 256k.

$\Diamond$ \textit{Applying LongRoPE search to broaden the fine-tuned extended LLM context window to 2048k.} As a concluding step, a secondary search is conducted using the LLM previously fine-tuned to a 256k context length. This process leads to an impressive increase in the context window size to 2048k, accomplished without additional rounds of fine-tuning. The overall achieved extension ratio amounts to $512\times$.

\textbf{Restoration in reduced context windows}. When the context window is greatly enlarged to 2048k, there is observable degradation in model performance within the initial context span. This phenomenon is documented with positional interpolation~\cite{pi}, which compresses the position embeddings for the original context into a restricted area of the embedding space, thereby impairing the efficacy of the language model. Employing a 512$\times$ enlargement results in substantial crowding of positional embeddings within the first 4k tokens, intensifying this challenge.

To address this issue, we conduct an additional evolutionary search on the expanded LLM, specifically to fine-tune the RoPE rescale parameters for shorter context spans such as 4k and 8k tokens. The search range for the maximal $\lambda$ is deliberately constrained, reflecting the decreased necessity for positional interpolation with limited context lengths. As the LLM processes input during inference, it adapts the relevant RoPE rescale factors in real-time.

\section{LongRoPE}

\label{sec:method}
Inspired by these observations, we propose LongRoPE. This approach begins by incorporating a high-performance search algorithm designed to leverage both types of non-uniformities, and subsequently applies it to expand the context window of LLMs past the 2 million token threshold.

\subsection{Problem Formulation}

These dual irregularities expand the solution landscape considerably and add layers of complexity to the optimization process. To manage these challenges, we recast the task of optimizing multidimensional non-uniform position interpolation as a search problem.

For a LLM targeting a  context window size of $L'$ and lengthy input documents $\mathbf{X}$, where each $\mathbf{x}\in\mathbf{X}$ surpasses $L'$ in token length, we denote the original rotary angle of the $i^{th}$ dimension in RoPE embedding at token position $n$ as  $\frac{n}{\beta^i}$. The optimization problem is then formulated as follows:

\begin{equation}
\begin{gathered}
		\label{eq:problem}

\underset {\mathbf{x}\in\mathbf{X}; \, |\mathbf{x}|\ge L'}{ \operatorname {arg\,min} } \,  \mathcal{L}\, \left( \text{LLM} (\text{RoPE}, \mathbf{X})\right), \text{with}
\\
\scalebox{0.93}{
$\underset {\begin{subarray}{l} i=0,\cdot\cdot,\frac{d}{2}-1;\\ n\in[ 0,|\mathbf{x}|);\end{subarray}}{\text{RoPE}_(n)}= \left[ \cdots, \cos\left( \mathbb{I}(\hat{\lambda}_{i}, \hat{n}) \cdot \frac{n}{\beta^i} \right), \sin\left( \mathbb{I}(\hat{\lambda}_{i}, \hat{n}) \cdot \frac{n}{\beta^i} \right), \cdots \right] $}\\
	\text{such that} \mathbb{I}(\hat{\lambda}_{i},\hat{n})=\begin{cases}
	1&\mbox{\;  $n<\hat{n}$}\\
{\frac{1}{\lambda}_{i}}&\mbox{\; $n\geq\hat{n}$}
\end{cases}
	\end{gathered}
\end{equation}

We define a group of rescale factors $\mathbb{I}(\hat{\lambda}_{i},\hat{n})$ to address both types of non-uniformities encountered. Here, $\hat{\lambda}_i$ signifies the irregularity in the RoPE dimension, while $\hat{n}$ represents the unevenness of token positions. The factor $\mathbb{I}(\hat{\lambda}_{i},\hat{n})$ adjusts the rotation angle for the $i^{th}$ RoPE dimension; specifically, $\hat{\lambda}_{i}$ acts as the scaling parameter and $\hat{n}$ sets the position cutoff. For the initial $\hat{n}-1$ token indices, the rotation retains the default RoPE formulation, namely $\frac{n}{\beta^i}$, without modification from the rescale factor. Once token positions reach $n \ge \hat{n}$, the rescale factor begins to influence the rotation angle.

Assuming a desired context window length of $L'$, the task is to determine the optimal rescale factors ($\mathbb{I}(\hat{\lambda}_{0},\hat{n})$, $\mathbb{I}(\hat{\lambda}_{1},\hat{n})$, ..., $\mathbb{I}(\hat{\lambda}_{i},\hat{n})$, ...) corresponding to the $1^{st}$ through $d^{th}$ RoPE dimensions. Through this approach, the adapted $\text{LLM}$ equipped with the rescaled $\text{RoPE}$ aims to minimize the next-token prediction loss, $\mathcal{L}$ (representing perplexity), for input batches $\mathbf{X}$ whose token sequence length is $L'$.

\subsection{Searching the Non-uniform Position Interpolation}

In addressing the issue outlined in Eq.~\ref{eq:problem}, we propose a straightforward and efficient approach that identifies the optimal RoPE rescale factors, enabling comprehensive utilization of the position embedding’s multidimensional non-uniform properties.

\textbf{Search space}. Our approach constructs an extensive search space that accommodates both forms of non-uniformity. Table~\ref{tbl:searchspace} details this search space configuration. In particular, we permit the assignment of a distinct rescale factor to each dimension within the RoPE embedding, enhancing flexibility. For ease in search space formulation, we elect to optimize over $\lambda_i$ and $\hat{n}$, rather than directly exploring $\mathbb{I}(\hat{\lambda}_{i},\hat{n})$, where $\hat{\lambda}_{i}=1/\lambda_i$. 
Table~\ref{tbl:searchspace} specifies that $\lambda_i$ ranges from 1.0 (corresponding to pure extrapolation) up to $s\times1.25$ (representing greater interpolation compared to PI), incremented by 0.01. Here, $s$ denotes the proportion by which the target context window is extended.

The parameter $\hat{n}$ determines how many of the initial token positions retain their original RoPE embeddings, without applying position interpolation. In practice, $\hat{n}$ is chosen from the set \{0, 1, 2, 4, 8, 12, 16, 20, 24, 28, 32, 64, 128, 256\}. If $\hat{n}=0$, then every token position relies exclusively on the searched rescale factors.

\textbf{Evolution-based search}. The search space described in Table~\ref{tbl:searchspace} includes a vast array of positional interpolation options, making it exceptionally difficult to traverse efficiently. For instance, a $s=4\times$ extension yields $400^{128/2}\times14$=4$\times10^{167}$ possible configurations. As the extension ratio increases, the search space expands at an exponential rate. To efficiently navigate this space, we employ evolution search~\cite{guo2020single} and incorporate two additional optimization strategies, significantly enhancing search performance. The detailed procedure is presented in Algorithm~\ref{alg:search}.

\textit{Optimized initial population generation}. Rather than generating all $P$ rescale factors at random to form the initial population, we explicitly include the three RoPE rescale factors associated with PI, NTK, and YaRN among the initial individuals. The rest of the $P$-3 individuals are produced by applying random mutations to these three rescale factors, each with a mutation probability $p$.

\textit{Monotonically non-decreasing constraint}. Once the initial population is produced, we evaluate LLM perplexity for every member. This entails setting the appropriate RoPE rescale factors within the target LLM and assessing the perplexity for input $\mathbf{X}$. The highest-scoring $k$ individuals are selected as parents in the evolutionary process. Nonetheless, the extensive search space allows simple mutation and crossover operations to wander toward suboptimal configurations, resulting in excessive and redundant perplexity evaluations. This inefficiency becomes pronounced, especially when $L'$ is large, due to the high computational cost associated with each perplexity inference.

To mitigate this issue, we enforce a monotonicity constraint whereby the rescaled RoPE factors are sampled such that $\lambda_i \leq \lambda_{i+1}$ for each index $i$. Only those RoPE configurations adhering to this constraint are utilized in LLM perplexity assessments, which notably streamlines the search process. In particular, we stipulate that $\lambda_i$ follows an increasing trend with respect to RoPE dimension, so that for indices $i=0,\ldots,63$, each successive factor is at least as large as its predecessor. This requirement of monotonicity across dimensions draws its rationale from NTK theory~\cite{jacot2018neural,tancik2020fourier,ntk}: it implies that dimensions associated with higher frequency (i.e., lower index values) necessitate less interpolation, yielding smaller values for $\lambda_i$, whereas dimensions corresponding to lower frequency (higher indices) accommodate greater interpolation, and thus permit larger $\lambda_i$.

\begin{algorithm}[t]
	\small
	\caption{The search algorithm}
	\textbf{Input:} target LLM, input samples $\mathbf{X}$,   population size $P$, mutation size $N_1$, crossover size $N_2$, max iterations $\mathcal{T}$, mutate probability $p$\\

	\begin{algorithmic}[1]
		\label{alg:search}
		\STATE Initialize Top-k as empty;	
		\STATE Set $\text{P}_0$ to \textit{Initialize\_population\_with\_optimization} ($P$, $\mathbf{X}$, $p$); 
		\FOR{$i$ from $1$ to $\mathcal{T}$}
		\STATE Apply \textit{Compute\_perplexity} (LLM, $\text{P}_{i-1}$, $\mathbf{X}$);
		\STATE  Update Top-k using \textit{Update\_Topk} (Top-k, $\text{P}_{i-1}$);
		\STATE Generate $\text{P}_{mutation}$ via \textit{Mutation\_with\_mono\_constraint} (Top-k, $N_1$, $p$); 
		\STATE Generate $\text{P}_{crossover}$ via \textit{Crossover\_with\_mono\_constraint} (Top-k, $N_2$);
		\STATE Construct $\text{P}_i$ as the union of $\text{P}_{mutation}$, $\text{P}_{crossover}$, and Top-k;
		\ENDFOR
		\STATE Return the Top-k member with the minimum perplexity;
	\end{algorithmic}
\end{algorithm}

\textbf{8$\times$ expansion without additional training}. The evolutionary search procedure efficiently discovers optimal non-uniform RoPE rescaling parameters, maintaining important dimensional features and positional encoding to reduce the degradation of information during interpolation. As illustrated in Fig.\ref{fig:non-ft}, our approach enables LLaMA2’s context window to grow from 4k to 32k in the absence of any further fine-tuning. By comparison, current techniques like PI as well as non-uniform NTK and YaRN result in a marked increase in perplexity beyond a 2$\times$ extension.

\subsection{Extending LLM Context Window to 2048K}
\label{sec:extendto2048k}

\textbf{Progressive extension to 2048k}. Here, we present an approach to expand the context window of pre-trained LLMs from the standard 4k tokens up to 2048k. Our experiments reveal that non-uniform positional interpolation enables an extension by a factor of 8$\times$ without the need for fine-tuning. However, achieving larger scaling (such as 512$\times$) necessitates fine-tuning. One possible strategy involves determining appropriate RoPE rescaling parameters at the target 2048k context size prior to fine-tuning. Yet, this approach is hampered by the extremely high computational costs involved in training. Furthermore, our observations indicate that fine-tuning LLMs under such large context extensions proves to be difficult in practice (refer to the Appendix).

LongRoPE demonstrates robust performance with both base and extended LLMs following fine-tuning. Accordingly, we propose a resource-efficient, incremental approach that attains the desired 2048k context window using only 1k fine-tuning iterations at a training sequence length of 256k.

$\Diamond$ \textit{Investigating the extension of pre-trained LLMs to a 256k window using LongRoPE search}. Using LLaMA2 as a case study, we search for optimal configurations aimed at target context lengths of 128k and 256k. At these settings, the expansion factor is 32$\times$ and 64$\times$, accordingly.

$\Diamond$ \textit{Fine-tuning to 256k}. Subsequently, we adapt the pre-trained LLM to support a context window of 256k through a staged fine-tuning approach. Initially, we perform 400 fine-tuning steps on LLaMA2 utilizing RoPE rescaled factors set for 128k. After this phase, we switch the RoPE rescaled factors to those corresponding to 256k and continue fine-tuning for another 600 steps. Compared to a straightforward fine-tuning to 256k, this two-step process demonstrates greater efficiency.

$\Diamond$ \textit{Scaling a previously fine-tuned extended LLM to 2048k via LongRoPE search}. In the final step, we apply an additional search procedure to the LLM already fine-tuned on a 256k context length. This enables us to reach an exceptionally large context window of 2048k, achieved without the need for any further fine-tuning. The overall extension is $512\times$ the original length.

\textbf{Recovering performance within shorter context windows}.  When increasing the context window size up to a substantial 2048k tokens, we observe diminished performance in the initial, smaller window range. This phenomenon has been documented in the literature on positional interpolation~\cite{pi}, which highlights that this technique compresses the positional embeddings at higher dimensions into a limited spatial region within the original window, thereby degrading the model's effectiveness. Specifically, extending the window by a factor of 512 results in the positions within the original 4k token window being densely packed together.

To address this issue, we conduct an additional evolution search on the expanded LLM, specifically tuning the RoPE rescale factors for shorter context windows (such as 4k and 8k tokens). For these shorter contexts, we restrict the upper bound of the searched $\lambda$ values, as less positional interpolation is needed. At inference time, the LLM adaptively modifies the relevant RoPE rescale factors in response to the input.

 \section{Experiments}

\label{sec:exp}
\subsection{Setup}
\textbf{Evaluation Tasks and Models}. We deploy LongRoPE with LLaMA2-7B and Mistral-7B, and assess their performance across three dimensions: (1) measuring perplexity for extended-context LLMs on lengthy texts; (2) evaluating performance on the Passkey retrieval task, which tests whether the model can extract a specific passkey amid a large amount of extraneous information; and (3) benchmarking against common LLM evaluation datasets restricted to a 4096-token context window.

\textbf{Fine-tuning}. For LLaMA2, fine-tuning is performed using a learning rate of 2e-5 with linear decay, employing a global batch size of 32. The model undergoes 400 training steps on the Redpajama~\cite{redpajama} dataset, which is divided into 128k-token segments and bounded by BOS and EOS tokens. Subsequently, from the resulting checkpoint, training continues for an additional 600 steps to extend the context window to 256k tokens. The 128k-context version is trained across 8 A100 GPUs with the distributed training framework~\cite{cube}; scaling up to 256k context necessitates 16 A100 GPUs. For Mistral, a constant learning rate of 1e-6 is selected alongside a global batch size of 64. Adhering to the YaRN~\cite{yarn} protocol, both 128k and 256k variants are trained for 400 steps using Together Computer's Long-Data Collections~\cite{mistraldata} with 16k sequence length, utilizing 4 A100 GPUs for the process.

\textbf{Search}. For target window sizes up to 256k, the following parameters are set: $P$=64, $N_1$=$N_2$=16, $p$=0.3, $\mathcal{T}$=40. In each generation, the top-32 candidates are chosen for the mutation/crossover steps. Perplexity assessment is conducted over 5 randomly selected samples from the PG19 validation set, ensuring each meets the minimum target context length. For window sizes exceeding 512k, we reduce the population size, as well as mutation and crossover pool sizes, by half. In this case, perplexity evaluation is based on 3 random examples drawn from the Pile-Books3~\cite{pile} validation set.

\textbf{Baselines}. To achieve a context window of 2048k, we performed fine-tuning on models initially trained with 128k and 256k context lengths. This results in LongRoPE-2048k (ft=128k) for LLaMA2 and LongRoPE-2048k (ft=256k) for Mistral. We evaluate these four models against leading context window extension baselines, focusing on open-source LLMs that have undergone fine-tuning following positional interpolation strategies such as PI, NTK, and YaRN. The comparison includes models like Together-32k~\cite{together}, Code LLaMA~\cite{codellama}, LongLoRA-full-FT-100k~\cite{longlora}, as well as YaRN-LLaMA and YaRN-Mistral~\cite{yarn}.

\subsection{Main Results}

\label{sec:mainresult}
\textbf{Language modeling on long sequences up to 256k}. Our analysis begins with a comparative study against the leading extended LLMs, all assessed within a 256k token evaluation window. To validate the robustness and adaptability of our approach, we select two distinct benchmarks: the test splits of Proof-pile~\cite{pg19} and PG19~\cite{pile}. Perplexity is measured at different context window sizes, leveraging a sliding window of size 256. For the PG19 corpus, all 100 documents comprising the full test set are used for evaluation. Regarding Proof-pile, aligning with YaRN~\cite{yarn}, we randomly draw 10 samples, each no shorter than 128k tokens.

Table~\ref{tbl:proof-pile} and Table~\ref{tbl:pg19} present a comparison of the perplexity scores for LLaMA2 and Mistral models, each augmented using various interpolation techniques, on the Proof-pile and PG19 datasets, respectively. Two principal findings are noteworthy: \textbf{(1)} our interpolated models exhibit a consistent decline in perplexity as evaluation lengths increase from 4k to 256k, demonstrating their enhanced capacity for utilizing extended context. \textbf{(2)} When tested with context windows up to 16$\times$ greater than typical settings—an environment known to impair performance on shorter spans—our LongRoPE-2048k variants still exceed the performance of leading baselines at the 256k context length.

\textbf{Language modeling for sequences exceeding 2000k tokens}. In order to assess performance on very large-scale documents, the Books3~\cite{pile} corpus serves as our testbed. To streamline evaluation, we randomly sample 20 books from the dataset, each with a length greater than 2048k tokens, and apply a sliding window of 256k tokens for processing.

As indicated in Table~\ref{tbl:books3}, LongRoPE effectively expands the context length for both LLaMA2-7B and Mistral-7B models up to 2048k, while maintaining perplexity metrics that are equal to or outperforming baseline values in shorter ranges from 8k to 128k. The results reveal distinct performance characteristics between LLaMA2 at 2048k and Mistral. While Mistral achieves superior results compared to baselines at shorter context spans, its perplexity rises above 7 for contexts exceeding 256k. The results for LLaMA2 generally conform to expectations: its perplexity improves as context size increases, with only minimal increases seen at the 1024k and 2048k ranges. Additionally, LongRoPE-2048k shows improved outcomes on LLaMA2 when the fine-tuning length is set to 256k versus 128k, attributable to a lower secondary extension ratio (namely, 8$\times$ rather than 16$\times$). On the other hand, Mistral yields better results when fine-tuned at a window size of 128k. This pattern primarily arises because, for both Mistral's 128k and 256k fine-tuning settings, training length is limited to 16k in accordance with YaRN's protocol, thus constraining Mistral’s capacity for context extension post fine-tuning.

\textbf{Passkey retrieval}. Next, we investigate the functional limits of context window size in generative tasks by employing a synthetic passkey retrieval evaluation, as introduced by ~\cite{passkey}. In this experiment, the model must locate and extract a random five-digit passkey embedded within an extensive document. The exact prompt format is outlined in the appendix. We conduct ten runs of the passkey retrieval procedure, ensuring that the passkey is positioned randomly with a uniform distribution spanning the entire context length used for evaluation.

Fig.~\ref{fig:passkey} presents a comparison of retrieval accuracy across different baseline approaches. Most existing models exhibit a steep decline in accuracy to zero once the token length surpasses 128k. Conversely, LongRoPE-LLaMA2-2048k (ft=256k) demonstrates strong performance, maintaining retrieval accuracy above 90\% for token ranges from 4k up to 2048k, even under the demanding scenario of extracting a passkey amidst millions of tokens. Similarly, LongRoPE-Mistral-2048k (ft=128k) sustains perfect (100\%) accuracy through 1800k tokens before falling to 60\% at 2048k. This trend is consistent with observations in Table~\ref{tbl:books3}, which shows a moderate increase in perplexity at the 2048k token level.

\textbf{Evaluation using established benchmarks in the initial context window}. The performance of LongRoPE-2048k models is assessed within their native context window by employing datasets from the Hugging Face Open LLM Leaderboard~\cite{huggingfaceboard} across both zero-shot and few-shot scenarios. Specifically, we utilize ARC-Challenge~\cite{allenai:arc} with 25 shots, HellaSwag~\cite{zellers2019hellaswag} with 10 shots, MMLU~\cite{mmlu} with 5 shots, and TruthfulQA~\cite{lin2021truthfulqa} in the zero-shot configuration.

According to Table~\ref{tbl:commonsense}, the models presented in this work yield results similar to those attained on the original benchmark, which was created for models with a more limited context window. In addition, they exceed the performance of the initial Mistral model on TruthfulQA by a margin of +0.5\%. Although LongRoPE-LLaMA2-2048k, which underwent fine-tuning at 256k tokens, displays marginally greater declines in performance, its results remain satisfactory across the majority of tasks.

\subsection{Ablation Results}

\textbf{Evaluating the impact of the second positional interpolation}. Within the progressive extension framework, we employ our search algorithm to apply a secondary non-uniform positional interpolation to the extended, fine-tuned LLMs. To assess the utility of this approach, we carry out experiments utilizing our fine-tuned LLaMA2-256k model, expanding its context window to 512k, 1024k, and 2048k tokens via PI and YaRN methods. As evidenced in Table~\ref{tbl:secondsearch}, our approach to non-uniform positional interpolation maintains perplexity at a relatively steady level. Conversely, perplexity associated with both PI and YaRN exhibits a sharp increase as the extension ratio grows.

\textbf{Recovery performance at reduced context lengths.} To address diminished effectiveness with shorter context windows, we fine-tune the RoPE scaling parameters for LongRoPE-2048k using our search methodology. In particular, we impose stricter limits on the maximum scaling factors during the search process to minimize interpolation for contexts of 4k and 8k. 
 Table~\ref{tbl:shortperformance} presents perplexity measurements for LongRoPE-LLaMA2-2048k evaluated on Proof-pile at 4k and 8k contexts, as well as aggregated LLM benchmark accuracy. The findings indicate a marked boost in performance at these shorter context sizes.

\textbf{Evaluation of the two types of non-uniformities}. To assess the individual impact of each non-uniformity, we conduct ablation studies. Two sets of experiments are performed: (i) LLaMA2-7B is extended to short contexts of 16k and 32k using three approaches—PI, searching solely across RoPE dimensions, and searching both non-uniformities; (ii) our fine-tuned LLaMA2 with context length 256k is scaled up to 2048k using the same methodologies. Perplexity is measured without any additional fine-tuning. According to Table~\ref{tbl:searchdimension}, incorporating non-uniformity in the RoPE dimension leads to a notable reduction in perplexity relative to linear interpolation via PI. Introducing non-uniformity in token positions enhances results at the shorter 16k and 32k sequences, but this benefit diminishes at 2048k, likely owing to the extremely lengthy context. Simply maintaining the original tokens without interpolation loses effectiveness, which we propose to further investigate in subsequent studies.

\section{Related Works}

Besides strategies centered on position interpolation, this section also reviews alternative methods proposed in related work.

\textbf{Retrieval-based} methods incorporate an external memory component to store information from previous contexts, while specialized retrieval modules are employed to obtain relevant documents during inference~\cite{longllama,wang2023augmenting,borgeaud2022improving}. These approaches often involve significant changes to the underlying LLM architectures. By contrast, our proposed method makes only modest adjustments to positional embeddings, resulting in a more streamlined solution. Additionally, our technique accommodates a broader range of long context applications, including long-form document summarization and few-shot learning scenarios, which are not limited to retrieval tasks.

\textbf{Attention-based context window extensions}. In addition to leveraging positional embedding interpolation, certain studies extend input context by maintaining the initial context window size of LLMs and altering the attention mechanism~\cite{han2023lminfinite, streamingllm,parallelwindow}. Central to these approaches is the reduction of the attention explosion problem associated with newly introduced positions, achieved through innovative attention masks. These techniques can be used in conjunction with positional interpolation approaches, as they serve complementary purposes.

\textbf{Fine-tuning based approaches} address the challenge of adapting pre-trained LLMs with revised position embeddings to handle extended contexts more efficiently. Studies such as Code LLaMA~\cite{codellama}, LLaMA2 Long~\cite{xiong2023effective}, and ScaledRoPE~\cite{liu2023scaling} implement a substantially increased base value for RoPE, subsequently fine-tuning models on the desired sequence length. In contrast, our method is compatible with diverse target lengths and is capable of handling sequences exceeding 2M in length. Furthermore, recent advancements including LongLoRA~\cite{longlora} and PoSE~\cite{zhu2023pose} have been introduced to alleviate the heavy GPU demands associated with fine-tuning models for extremely long contexts (i.e., over 128k tokens). Importantly, our approach remains complementary to these efficient fine-tuning strategies.

\section{Conclusion}

This paper introduces LongRoPE, a technique that dramatically increases the context window of LLMs to 2048k tokens, surpassing previous limits, while preserving their effectiveness with shorter sequences. Our approach leverages two distinct types of non-uniformities in RoPE positional embeddings, utilizing a computationally efficient evolutionary search. This not only ensures robust initial conditions for subsequent fine-tuning but also allows for an 8$\times$ increase in context window size without requiring fine-tuning. Furthermore, we advocate a gradual expansion methodology, employing LLMs fine-tuned on 256k contexts to ultimately achieve a 2048k window without additional fine-tuning steps. Comprehensive experimental analysis substantiates LongRoPE’s performance. We anticipate that LongRoPE-2048k will support a wide array of long context applications and foster new directions in this domain.

\section*{Broader Impacts} The primary objective of this research is to contribute to the advancement of Machine Learning as a discipline. While our study may have various implications for society, we have not identified any particular issues that require special emphasis at this time.

	\balance

\end{document}